\documentclass[12pt]{article}
\usepackage[utf8]{inputenc}
\usepackage[T1]{fontenc}
\usepackage{hyperref}
\title{Taliabu Language Resources for K-12 Curriculum Development}
\author{Compiled by Claude Sonnet 5 (claude-sonnet-5) for Andr\'as Kornai}
\date{}
\begin{document}
\maketitle

\section{Introduction}
Taliabu (also Taliabo; ISO 639-3 tlv; Glottocode tali1262) is an Austronesian language of Taliabu Island in the Sula Islands, North Maluku, Indonesia, with roughly 4,900 speakers. The dialects named in the sources found are Kadai, Padang (Samala), Mananga and Mangei (Soboyo). Its subgrouping is unsettled. Searches found no school materials, primers, textbooks or a published dictionary for Taliabu. This dossier is therefore short.

\section{Primary Educational Resources}
No primers, textbooks or school programmes for Taliabu turned up in the search results. Nothing is known to exist.

\section{Linguistic and Reference Materials}
\begin{itemize}
\item Wikipedia, ``Taliabo language'': an overview of speakers, dialects and classification. \url{https://en.wikipedia.org/wiki/Taliabo_language}
\item ``Phonological Variation of Taliabu Language Dialects'' (journal article, Retorika): compares the Kadai and Mange dialects. \url{https://ojs.unm.ac.id/retorika/article/view/9266/0}
\item ``Sula: Its language, land, and people'' (University of Hawai\textquoteleft i dissertation, 2020): a grammatical sketch and lexicon of Sula, which is a different language (ISO szn). It covers a dialect continuum that reaches Taliabu island, so it is useful for comparison only. \url{https://scholarspace.manoa.hawaii.edu/items/daeb65e2-8b98-4002-93b2-ced62f91fe4d}
\item Wikipedia, ``Sula language'': background on the neighbouring language. \url{https://en.wikipedia.org/wiki/Sula_language}
\end{itemize}
The search summary also mentioned a 1921 Dutch description by Fortgens of the Sobojo variety. It was not located online here, so no URL is given.

\section{Religious and Community Resources}
No scripture translations or community organisations for Taliabu were found in the search results.

\section{Digital and Audio Resources}
No audio, video, corpus or software resources were found.

\section{Curriculum Development Context}
Taliabu appears to be undocumented for educational purposes. A curriculum would have to start from new documentation: an agreed orthography, a lexicon and recorded texts, ideally made with the speaker community. The published phonology and morphology articles and the Sula work are the only linguistic starting points found.
\end{document}